%! Inverse Functions — Unit 6, Lesson 6.7 — Group Activity (Answer Key)
\documentclass[10pt]{article}
\usepackage{algebra2-article}
\usepackage{algebra2-key}   % pulls in -boxes; do NOT also load -boxes
\newcommand{\TallMath}[1]{$\displaystyle #1\rule[-1.4em]{0pt}{3.2em}$}
% In-formula answer slot (\ans is text-mode and must never appear inside $...$).
\newcommand{\mans}[1]{{\color{keyred}\mathbf{#1}}}
\newcommand{\mexp}[1]{{\color{keyred}#1}}
\newcommand{\lgrid}[4]{%
  \draw[step=1, linegray!40, very thin] (#1,#3) grid (#2,#4);
  \draw[->, semithick, charcoal] (#1,0)--(#2,0) node[below right, font=\scriptsize]{$x$};
  \draw[->, semithick, charcoal] (0,#3)--(0,#4) node[left, font=\scriptsize]{$y$};}

\begin{document}

\pageheader{Unit 6, Lesson 6.7}{Group Activity --- Answer Key}
\namepartnerperiod

\vspace{0.08in}

\begin{headlinebox}{forestbg}
\small\textbf{Undo It.} Every function today comes with a partner whose only job is to put things
back the way they were. Two rules for the period: \textbf{(1)} reverse the operations \emph{and} the
order, and \textbf{(2)} never announce an inverse you have not tested by composing \emph{both} ways.
\end{headlinebox}

\vspace{0.06in}

\begin{tcolorbox}[colback=white, colframe=black!40, title=\textbf{Tier R --- Remediate},
    fonttitle=\bfseries, arc=1.5mm, left=3mm, right=3mm, top=2mm, bottom=2mm, breakable]
\textbf{Part 1 --- Undo the machine.} List what the function does, in order. Then list the undoing
steps \emph{in the reverse order}, and write the rule.
\begin{center}
\renewcommand{\arraystretch}{1.8}
\begin{tabularx}{\linewidth}{@{}l X X X@{}}
\hline
\textbf{$f(x)$} & \textbf{What it does, in order} & \textbf{Undo, in reverse order} & \textbf{$f^{-1}(x)$} \\
\hline
(a) $x+8$ & add $8$ & \ans{subtract $8$} & \ans{$x-8$} \\
(b) $3x$ & multiply by $3$ & \ans{divide by $3$} & \ans{$\frac{x}{3}$} \\
(c) $2x-1$ & multiply by $2$, then subtract $1$ & \ans{add $1$, then divide by $2$} & \ans{$\frac{x+1}{2}$} \\
(d) $\dfrac{x}{5}+4$ & \ans{divide by $5$, then add $4$} & \ans{subtract $4$, then multiply by $5$} & \ans{$5x-20$} \\
\hline
\end{tabularx}
\end{center}
\vspace{-4pt}
In rows (c) and (d) the undoing steps came out in the \emph{opposite} order from the original.
Explain to your partner why that has to happen. \ans{The last thing done is the first thing that has
to come off --- shoes before socks.}

\vspace{6pt}
\textbf{Part 2 --- Test the partner.} Compose both ways. Only one verdict per row, and only after
\emph{both} boxes are filled.
\begin{center}
\renewcommand{\arraystretch}{1.9}
\begin{tabularx}{\linewidth}{@{}l X X c@{}}
\hline
\textbf{$f$ and $g$} & \textbf{$f\big(g(x)\big)$} & \textbf{$g\big(f(x)\big)$} & \textbf{Inverses?} \\
\hline
(a) $f(x)=x+7$, \; $g(x)=x-7$ & \ans{$(x-7)+7=x$} & \ans{$(x+7)-7=x$} & \ans{yes} \\
(b) $f(x)=4x$, \; $g(x)=\dfrac{x}{4}$ & \ans{$4\cdot\frac{x}{4}=x$} & \ans{$\frac{4x}{4}=x$} & \ans{yes} \\
(c) $f(x)=3x+1$, \; $g(x)=\dfrac{x}{3}+1$ & \ans{$3\!\left(\frac{x}{3}+1\right)+1=x+4$} & \ans{$\frac{3x+1}{3}+1=x+\frac{4}{3}$} & \ans{no} \\
\hline
\end{tabularx}
\end{center}
\vspace{0.25in}
Row (c) uses the right two operations but gets them in the wrong order. Fix it: the real inverse of
$f(x)=3x+1$ is $f^{-1}(x)=$ \ans{$\dfrac{x-1}{3}$}

\vspace{6pt}
\textbf{Part 3 --- Swap the points.} \; Here is a table of $f$, and the graph of a different
function $p$.
\begin{center}
\begin{minipage}[c]{0.50\linewidth}
\centering\small
\renewcommand{\arraystretch}{1.4}
\begin{tabular}{@{}|c|c|c|c|c|c|@{}}
\hline
$x$ & $-2$ & $0$ & $1$ & $3$ & $5$ \\
\hline
$f(x)$ & $-5$ & $-1$ & $1$ & $5$ & $9$ \\
\hline
\end{tabular}
\\[8pt]
\renewcommand{\arraystretch}{1.4}
\begin{tabular}{@{}|c|c|c|c|c|c|@{}}
\hline
$x$ & \ans{$-5$} & \ans{$-1$} & \ans{$1$} & \ans{$5$} & \ans{$9$} \\
\hline
$f^{-1}(x)$ & \ans{$-2$} & \ans{$0$} & \ans{$1$} & \ans{$3$} & \ans{$5$} \\
\hline
\end{tabular}
\end{minipage}
\hfill
\begin{minipage}[c]{0.46\linewidth}
\centering
\begin{tikzpicture}[scale=0.36]
  \lgrid{-5}{5}{-5}{5}
  \draw[dashed, semithick, charcoal] (-4.6,-4.6)--(4.6,4.6);
  \draw[very thick, forest] (-4,-3)--(0,1)--(4,3);
  \fill[forest] (-4,-3) circle (3.4pt);
  \fill[forest] (0,1) circle (3.4pt);
  \fill[forest] (2,2) circle (3.4pt);
  \fill[forest] (4,3) circle (3.4pt);
  \node[forest, font=\scriptsize] at (-2.4,2.6) {$y=p(x)$};
\end{tikzpicture}
\end{minipage}
\end{center}
\vspace{-2pt}
Four points are marked on the graph of $p$: $(-4,-3)$, $(0,1)$, $(2,2)$, $(4,3)$. Write the four
matching points on the graph of $p^{-1}$.
\begin{center}
\ans{$(-3,-4)$} \hspace{0.2in} \ans{$(1,0)$} \hspace{0.2in} \ans{$(2,2)$} \hspace{0.2in}
\ans{$(3,4)$}
\end{center}
\vspace{-4pt}
One of those four points did not move at all. Which one, and what is special about where it sits?
\ans{$(2,2)$ --- it lies on the mirror line $y=x$.}
\end{tcolorbox}

\begin{tcolorbox}[colback=white, colframe=black!40, title=\textbf{Tier A --- Approaching Proficiency},
    fonttitle=\bfseries, arc=1.5mm, left=3mm, right=3mm, top=2mm, bottom=2mm, breakable]
\textbf{Part 1 --- Swap and solve.} Two of these six inverses need a restriction. Find them.
\begin{center}
\renewcommand{\arraystretch}{1.9}
\begin{tabularx}{\linewidth}{@{}l X l X@{}}
\hline
\textbf{$f(x)$} & \textbf{$f^{-1}(x)$} & \textbf{$f(x)$} & \textbf{$f^{-1}(x)$} \\
\hline
(a) $4x+9$ & \ans{$\frac{x-9}{4}$} & (d) $(x+5)^{3}$ & \ans{$\sqrt[3]{x}-5$} \\
(b) $\dfrac{2x-1}{3}$ & \ans{$\frac{3x+1}{2}$} & (e) $\sqrt{x+5}$ & \ans{$x^{2}-5$, $x\ge0$} \\
(c) $x^{3}-4$ & \ans{$\sqrt[3]{x+4}$} & (f) $\sqrt{x}-3$ & \ans{$(x+3)^{2}$, $x\ge-3$} \\
\hline
\end{tabularx}
\end{center}
\vspace{2pt}
{\small \emph{Work, the two that bite:} \; (e) $x=\sqrt{y+5}\Rightarrow x^{2}=y+5\Rightarrow
y=x^{2}-5$, and $f$ only ever outputs numbers $\ge0$, so $f^{-1}$ may only accept $x\ge0$. \;
(f) $x=\sqrt{y}-3\Rightarrow x+3=\sqrt{y}\Rightarrow y=(x+3)^{2}$, and $\sqrt{y}\ge0$ forces
$x+3\ge0$, i.e.\ $x\ge-3$. \; Rows (c) and (d) need \emph{no} restriction --- an odd index accepts
every real number.}

\vspace{0.10in}

\textbf{Part 2 --- Justify it.} Pick row (a) and row (e). For each, show \emph{both} compositions and
state the verdict.
\ansline{(a) $f\!\left(\frac{x-9}{4}\right)=4\!\left(\frac{x-9}{4}\right)+9=x$ \; and \;
$f^{-1}(4x+9)=\frac{(4x+9)-9}{4}=x$. \; Inverses.}
\ansline{(e) $f(x^{2}-5)=\sqrt{(x^{2}-5)+5}=\sqrt{x^{2}}=|x|=x$ \emph{because} $x\ge0$ on the
domain of $f^{-1}$ --- this is where the restriction earns its keep.}
\ansline{\ \ \ \ And $f^{-1}\!\left(\sqrt{x+5}\,\right)=\left(\sqrt{x+5}\,\right)^{2}-5=(x+5)-5=x$
for every $x\ge-5$. \; Inverses, once the restriction is stated.}
\ansline{Verdict for both rows: yes. Neither is settled by one order alone --- row (e) is exactly the
$\left(\sqrt{x}\,\right)^{2}$ versus $\sqrt{x^{2}}$ pair from Lesson 6.6, in disguise.}

\vspace{4pt}
\textbf{Part 3 --- Does an inverse function exist?} Apply the horizontal line test to each
pre-drawn graph.
\begin{center}
\begin{minipage}[c]{0.235\linewidth}
\centering
\begin{tikzpicture}[scale=0.26]
  \lgrid{-5}{5}{-5}{5}
  \draw[very thick, forest] (-5,-1.5)--(5,3.5);
  \node[forest, font=\tiny] at (-4.0,-4.0) {(a)};
\end{tikzpicture}
\end{minipage}
\hfill
\begin{minipage}[c]{0.235\linewidth}
\centering
\begin{tikzpicture}[scale=0.26]
  \lgrid{-5}{5}{-5}{5}
  \begin{scope}
    \clip (-5,-5) rectangle (5,5);
    \draw[very thick, navy, domain=-3:3, samples=110, smooth] plot (\x,{(\x)*(\x)-4});
  \end{scope}
  \node[navy, font=\tiny] at (-4.0,-4.0) {(b)};
\end{tikzpicture}
\end{minipage}
\hfill
\begin{minipage}[c]{0.235\linewidth}
\centering
\begin{tikzpicture}[scale=0.26]
  \lgrid{-5}{5}{-5}{5}
  \begin{scope}
    \clip (-5,-5) rectangle (5,5);
    \draw[very thick, forest, domain=0:5, samples=110, smooth] plot (\x,{sqrt(\x)});
  \end{scope}
  \fill[forest] (0,0) circle (3pt);
  \node[forest, font=\tiny] at (-4.0,-4.0) {(c)};
\end{tikzpicture}
\end{minipage}
\hfill
\begin{minipage}[c]{0.235\linewidth}
\centering
\begin{tikzpicture}[scale=0.26]
  \lgrid{-5}{5}{-5}{5}
  \begin{scope}
    \clip (-5,-5) rectangle (5,5);
    \draw[very thick, navy, domain=-2.3:2.3, samples=140, smooth]
      plot (\x,{(\x)*(\x)*(\x)-3*(\x)});
  \end{scope}
  \node[navy, font=\tiny] at (-4.0,-4.0) {(d)};
\end{tikzpicture}
\end{minipage}
\end{center}
\vspace{-2pt}
\begin{center}
\renewcommand{\arraystretch}{1.7}
\begin{tabularx}{0.96\linewidth}{@{}l c X@{}}
\hline
 & \textbf{Passes the test?} & \textbf{If not: how many times does the \emph{worst} horizontal line cross?} \\
\hline
(a) & \ans{yes} & \ans{---\; a line never repeats an output} \\
(b) & \ans{no} & \ans{$2$ (every line above the vertex)} \\
(c) & \ans{yes} & \ans{---\; the arm only climbs} \\
(d) & \ans{no} & \ans{$3$ (e.g.\ $y=1$, between the two turning points)} \\
\hline
\end{tabularx}
\end{center}
\vspace{-4pt}
For graph (b), name a restricted domain that would let it have an inverse function.
\ans{$x\ge0$ (or $x\le0$) --- either arm alone}

\vspace{4pt}
\textbf{Part 4 --- Say what you found.}
\begin{enumerate}[label=\textbf{J\arabic*.}, leftmargin=2.4em, itemsep=6pt]
  \item Row (f) of Part 1 is the one everybody gets \emph{almost} right. Its inverse needs the
        condition $x\ge-3$, not $x\ge0$. Where does the $-3$ come from? \ansline{From the
        \emph{range} of $f$. Since $\sqrt{x}\ge0$, the outputs of $f(x)=\sqrt{x}-3$ are all
        $\ge-3$ --- and the outputs of $f$ are exactly the inputs $f^{-1}$ is allowed to take.}
        \ansline{The restriction is never guessed; it is read off the original function's range.}
  \item Complete the table for row (e), $f(x)=\sqrt{x+5}$, and say in one sentence what you notice.
    \begin{center}
    \renewcommand{\arraystretch}{1.6}
    \begin{tabularx}{0.9\linewidth}{@{}l X X@{}}
    \hline
     & \textbf{Domain} & \textbf{Range} \\
    \hline
    $f$ & \ans{$x\ge-5$} & \ans{$y\ge0$} \\
    $f^{-1}$ & \ans{$x\ge0$} & \ans{$y\ge-5$} \\
    \hline
    \end{tabularx}
    \end{center}
    \ansline{The two rows are the same pair of sets with the columns swapped: domain and range
    trade places.}
  \item A classmate writes: ``The inverse of $f(x)=x^{2}-4$ is $f^{-1}(x)=\sqrt{x+4}$.'' Their
        algebra is fine. Using your answer to Part 3(b), say what is wrong with the
        \emph{sentence}, and rewrite it so it is true. \ansline{$f(x)=x^{2}-4$ fails the horizontal
        line test on its full domain, so it has \emph{no} inverse function at all --- there is
        nothing for ``the inverse'' to name. (Check: $f(-3)=f(3)=5$, so $f^{-1}(5)$ would have to be
        both $-3$ and $3$.)}
        \ansline{True version: ``If $f(x)=x^{2}-4$ is restricted to $x\ge0$, then
        $f^{-1}(x)=\sqrt{x+4}$.''}
        \ansline{The algebra was never the problem. Swap-and-solve will happily produce a formula
        for a function that has no inverse; the horizontal line test is what tells you whether the
        formula is entitled to the name.}
  \item In Tier R Part 3, one marked point stayed exactly where it was. State the general rule:
        which points of a graph never move when you reflect over $y=x$, and why?
        \ansline{The points already \emph{on} the line $y=x$ --- the ones with $x=y$. Swapping the
        coordinates of $(2,2)$ gives $(2,2)$ back, so a reflection cannot move it.}
        \ansline{Consequence used in Tier E: if $f$ and $f^{-1}$ ever meet, they meet on $y=x$.}
\end{enumerate}
\end{tcolorbox}

\begin{tcolorbox}[colback=white, colframe=black!40, title=\textbf{Tier E --- Extension},
    fonttitle=\bfseries, arc=1.5mm, left=3mm, right=3mm, top=2mm, bottom=2mm, breakable]
\textbf{Part 1 --- Finish the Hook.} $F(c)=\dfrac{9}{5}c+32$ turns Celsius into Fahrenheit.
\begin{enumerate}[label=(\alph*), leftmargin=1.9em, itemsep=6pt]
  \item Build $F^{-1}$ by swapping and solving. \; $F^{-1}(f)=$ \ans{$\frac{5}{9}(f-32)$}

        {\small \emph{Work:} $c=\frac95 f'+32$ after the swap, so $c-32=\frac95 f'$ and
        $f'=\frac59(c-32)$ --- subtract \emph{first}, then scale. Malik's version.}
  \item Prove it is right --- both orders, algebraically, not just at one temperature.
        \ansline{$F\big(F^{-1}(f)\big)=\frac95\cdot\frac59(f-32)+32=(f-32)+32=f$ \checkmark}
        \ansline{$F^{-1}\big(F(c)\big)=\frac59\!\left(\left(\frac95 c+32\right)-32\right)
        =\frac59\cdot\frac95 c=c$ \checkmark \; --- Rosa's $\frac59f-32$ fails both.}
        \ansline{One test temperature could never settle this; two functions can agree at a point
        and disagree everywhere else.}
  \item There is exactly one temperature that reads the \emph{same} number on both scales. Find it
        by solving $F(c)=c$. \; $c=$ \ans{$-40$}

        {\small \emph{Work:} $\frac95 c+32=c \Rightarrow 32=c-\frac95 c=-\frac45 c \Rightarrow
        c=-40$. \; $-40^\circ$C $=-40^\circ$F.}
  \item Explain that answer using the reflection picture: where must the graphs of $F$ and $F^{-1}$
        cross, and why is that the same as asking for $F(c)=c$? \ansline{A graph and its reflection
        can only meet on the mirror line, so $F$ and $F^{-1}$ cross on $y=x$.}
        \ansline{A point on $y=x$ is a point whose output equals its input --- which is exactly the
        equation $F(c)=c$. Same question, asked geometrically.}
\end{enumerate}

\vspace{4pt}
\textbf{Part 2 --- Functions that are their own inverse.} Find $f^{-1}$ for each.
\begin{center}
\renewcommand{\arraystretch}{1.8}
\begin{tabularx}{\linewidth}{@{}X X X@{}}
(e) \; $f(x)=6-x$ \; $f^{-1}(x)=$ \ans{$6-x$} &
(f) \; $f(x)=\dfrac{1}{x}$ \; $f^{-1}(x)=$ \ans{$\dfrac{1}{x}$} &
(g) \; $f(x)=\dfrac{2x+3}{x-2}$ \; $f^{-1}(x)=$ \ans{$\dfrac{2x+3}{x-2}$} \\
\end{tabularx}
\end{center}
\vspace{2pt}
{\small \emph{Work for (g):} $x=\frac{2y+3}{y-2}\Rightarrow x(y-2)=2y+3\Rightarrow xy-2x=2y+3
\Rightarrow y(x-2)=2x+3\Rightarrow y=\frac{2x+3}{x-2}$ --- the same rule it started as.}
\begin{enumerate}[label=(\alph*), leftmargin=1.9em, itemsep=5pt, start=8]
  \item All three came out the same as $f$. What does that force to be true about each graph and the
        line $y=x$? \ansline{Each graph is its own reflection over $y=x$ --- it is symmetric about
        that line. (Reflecting it does nothing at all.)}
        \ansline{Such a function is called an \emph{involution}: applying it twice returns $x$.}
\end{enumerate}

\vspace{4pt}
\textbf{Part 3 --- Socks and shoes, in symbols.} Let $f(x)=2x$ and $g(x)=x+3$.
\begin{enumerate}[label=(\alph*), leftmargin=1.9em, itemsep=6pt, start=9]
  \item $(f\circ g)(x)=$ \ans{$2x+6$}, \; so $(f\circ g)^{-1}(x)=$ \ans{$\frac{x}{2}-3$}
  \item $f^{-1}(x)=$ \ans{$\frac{x}{2}$} and $g^{-1}(x)=$ \ans{$x-3$}. Now build \emph{both} orders: \\[3pt]
        $g^{-1}\big(f^{-1}(x)\big)=$ \ans{$\frac{x}{2}-3$ \checkmark} \hspace{0.3in}
        $f^{-1}\big(g^{-1}(x)\big)=$ \ans{$\frac{x-3}{2}$ \; $\times$}
  \item Only one of those two matched your answer to (i). Complete the rule:
        $(f\circ g)^{-1}=$ \ans{$g^{-1}\circ f^{-1}$}
  \item Say why in plain English, using socks and shoes. \ansline{Socks first, then shoes. To undo
        it you take the \emph{shoes} off first --- the machine that ran \textbf{last} is the one
        undone \textbf{first}, so the inverses come back in the opposite order.}
        \ansline{This is Warm-Up item 3 and the Hook, now proved in general.}
\end{enumerate}

\vspace{4pt}
\textbf{Part 4 --- The other half of the parabola.} In the notes you restricted $y=x^{2}$ to
$x\ge0$ and got $y=\sqrt{x}$.
\begin{enumerate}[label=(\alph*), leftmargin=1.9em, itemsep=5pt, start=13]
  \item Restrict it to $x\le0$ instead --- the \emph{left} arm. Swap and solve, and keep the arm you
        are actually on. \; $f^{-1}(x)=$ \ans{$-\sqrt{x}$}

        {\small \emph{Work:} $x=y^{2}$ with $y\le0$, so $y=-\sqrt{x}$ --- the negative root is the
        one that lives on the arm you kept.}
  \item Both restrictions are legal, and they give different inverses. So ``the'' inverse of
        $y=x^{2}$ is not a well-posed question until you say what? \ans{which half of the parabola
        you kept --- i.e.\ the restricted domain}
  \item The area of a circle is $A=\pi r^{2}$ --- the same shape of function. Solve it for $r$, and
        say which restriction the \emph{situation} chooses for you, without any algebra.
        \ansline{$r=\sqrt{A/\pi}$. The context restricts $r\ge0$ all by itself: a radius cannot be
        negative, so the left arm was never on the table.}
        \ansline{This is the same move as Lesson 6.0's free-fall and skid-mark models --- context
        picks the branch that algebra alone leaves open.}
\end{enumerate}
\end{tcolorbox}

\begin{teachernote}
\small \textbf{Tier R} is deliberately arithmetic-first: Part 1 never uses the word ``inverse'' until
the last column, because the reversal of \emph{order} is the idea that has to land, and it lands
faster on ``multiply by $2$, then subtract $1$'' than on notation. Row (c) of Part 2 is the whole
lesson in one line --- right operations, wrong order, and only composition catches it.
\\[3pt]
\textbf{Tier A} rows (e) and (f) of Part 1 are the assessed discriminator. Nearly everyone produces
$x^{2}-5$ and $(x+3)^{2}$; almost nobody writes the restriction unprompted, and the students who
write $x\ge0$ on row (f) have pattern-matched rather than thought. \textbf{J1 is the item to
circulate for} --- the $-3$ comes from the \emph{range} of $f$, and a student who can say that
sentence has the whole concept. J3 is the deepest question on the page: swap-and-solve will happily
hand you a formula for a function that has no inverse, so the horizontal line test is what licenses
the word ``the.''
\\[3pt]
\textbf{Tier E} Part 1 is worth previewing to the whole class even if only one group reaches it ---
$-40^\circ$ is a fact students remember for years, and part (d) explains it with the mirror line
rather than with algebra. Part 3 is the formal version of the Hook and of Warm-Up item 3: reversing
a composition reverses the order of the inverses, which is why Rosa's answer was wrong. If a group
finishes Part 3, send them to Part 4(o) and ask them to connect it to Lesson 6.0's free-fall model.
\end{teachernote}

\end{document}
